% !PRINTED 524
\weberpage{524}

%<<CONT
where for $\gamma_2(\omega)$ and $j(\omega)$ the values (5) are to be substituted. These equations can be decomposed still further, and for $x$ itself one obtains much simpler equations:
\weq{8}{\begin{array}{ll}
f(\sqrt{-51})^3 = 2x, & x^3 - 4x^2 - x - 1 = \sqrt{17}\,x^2,\\
f(\sqrt{-91}) = x, & x^3 - 2x^2 - x - 2 = \sqrt{13}\,x,\\
f(\sqrt{-99})^3 = 2x, & x^3 - 13x^2 - 4x - 1 = \sqrt{33}\left(2x^2 + x\right).
\end{array}}
The representation of the class invariants $j(\omega)$ by a single square root is always possible whenever two genera are present and one class in each genus. These values of $m$ are found in \gesperrt{Gauss}'s table of class numbers (Werke, Bd.\,II, S.\,450\,ff.) under the designation II, 3, and of these the ones to be selected are those which are $\equiv 3\ (\mathrm{mod}\ 8)$. Their number is in all probability finite, as the table just mentioned shows; they are the numbers:
\[m = 35, \quad 51, \quad 75, \quad 91, \quad 99, \quad 115, \quad 123,\]
\[147, \quad 187, \quad 235, \quad 267, \quad 403, \quad 427.\]
If the classes of forms split into an arbitrary number of genera, but each genus contains only one class, then by the theorem proved in § 138 all class invariants can be expressed by \gesperrt{square roots}. \gesperrt{Euler} and \gesperrt{Gauss} concluded by induction that there is only a finite number of such discriminants, namely 65\footnote{\gesperrt{Euler}, Nouv. Mém. de Berlin 1776, S.\,338. \gesperrt{Gauss}, Disq. art. 303. \gesperrt{Euler} arrived at these numbers by a different route, namely from the problem of determining large primes (Cf. the Straßburg dissertation of \gesperrt{Peter Meyer}: Beweis eines von \gesperrt{Euler} entdeckten Satzes, betreffend die Bestimmung von Primzahlen, Straßburg 1906). These numbers are:

1, 2, 3, 4, 5, 6, 7, 8, 9, 10, 12, 13, 15, 16, 18, 21, 22, 24, 25, 28, 30, 33, 37, 40, 42, 45, 48, 57, 58, 60, 70, 72, 78, 85, 88, 93, 102, 105, 112, 120, 130, 133, 165, 168, 177, 190, 210, 232, 240, 253, 273, 280, 312, 330, 345, 357, 385, 408, 462, 520, 760, 840, 1320, 1365, 1848.

% Source chapter-closing rule follows the completed footnote.
\par\medskip\noindent\makebox[\linewidth]{\rule{0.25\linewidth}{0.2pt}}}. In the memoir \glqq Zur komplexen Multiplikation elliptischer Funktionen\grqq{} (Mathematische Annalen, Bd.\,33) I have computed all these numbers. They are also contained in the table of class invariants which is appended to this book.
